\documentclass[article]{jss}

\usepackage{array}

\raggedbottom
\newcolumntype{L}[1]{>{\raggedright\arraybackslash}p{#1}}
\setlength{\textfloatsep}{10pt plus 2pt minus 3pt}
\setlength{\floatsep}{10pt plus 2pt minus 3pt}
\setlength{\intextsep}{10pt plus 2pt minus 3pt}
\setlength{\abovecaptionskip}{6pt plus 1pt minus 1pt}
\setlength{\belowcaptionskip}{2pt plus 1pt minus 1pt}
\emergencystretch=2em


%% -- Article metainformation -------------------------------------------------

\author{Matthew Powell\\Independent researcher}
\Plainauthor{Matthew Powell}

\title{CausalLens: Diagnostics-First Causal Inference Software in \proglang{Python}}
\Plaintitle{CausalLens: Diagnostics-First Causal Inference Software in Python}
\Shorttitle{CausalLens in \proglang{Python}}

\Abstract{
  CausalLens is a \proglang{Python} package for causal effect estimation,
  diagnostics, sensitivity analysis, and manuscript-oriented reporting across
  observational and quasi-experimental designs. The package combines regression
  adjustment, propensity-score matching, inverse probability weighting (IPW),
  doubly robust (DR) estimation, cross-fitted doubly robust estimation,
  difference-in-differences (DiD), synthetic control, instrumental variables (IV),
  regression discontinuity (RD), manipulation testing, and bunching analysis in a
  shared diagnostics-first workflow. The contribution is not a new estimator or
  identification strategy. Instead, CausalLens provides a compact software
  layer that keeps overlap checks, balance summaries, sensitivity diagnostics,
  reference-parity checks, and exportable outputs close to the effect estimate.
  This positions the project as review-oriented statistical software: it helps
  users fit established causal estimators while retaining the assumption and
  validation context needed for replication and peer review. In the current
  replication bundle, Lalonde benchmark estimates fall between \$1429 and
  \$1710 across four estimators, while NHEFS smoking-cessation estimates
  fall between 3.26 and 3.33~kg, illustrating that the package reproduces
  recognizable benchmark ranges while exporting the associated diagnostics.
}

\Keywords{causal inference, statistical software, diagnostics,
  reproducibility, treatment effects, validation, \proglang{Python}}
\Plainkeywords{causal inference, statistical software, diagnostics,
  reproducibility, treatment effects, validation, Python}

\Address{
  Matthew Powell\\
  Independent researcher\\
  E-mail: \email{mattpowell.gis@gmail.com}\\
  URL: \url{https://www.linkedin.com/in/matthew-a-powell/}
}

\begin{document}


%% -- Introduction ------------------------------------------------------------

\section[Introduction: Reviewable causal inference software in Python]{Introduction: Reviewable causal inference software in \proglang{Python}}\label{sec:intro}

The potential-outcomes framework introduced by \citet{Rubin1974} and
formalized within the ``Rubin Causal Model'' \citep{Holland1986} provides the
conceptual foundation for modern treatment-effect estimation. Under this
framework, every causal effect estimate arrives with a burden of proof:
overlap assumptions, covariate balance, sensitivity to hidden bias, and
enough documented context for another analyst to judge whether the result is
credible. Recent textbook treatments emphasize these requirements from both
econometric \citep{AngristPischke2009,Wooldridge2010,Cunningham2021} and
statistical perspectives
\citep{HernanRobins2020,HuntingtonKlein2022,Ding2024}. In practice, that
burden is distributed across separate tools, scripts, and post-hoc checks
that rarely travel with the estimate itself
\citep{Imbens2004,ImbensWooldridge2009,AtheyImbens2017,Peng2011,SandveEtAl2013,StoddenEtAl2016}.

The existing \proglang{Python} ecosystem reflects this fragmentation.
Table~\ref{tab:comparison} summarizes the landscape. \pkg{DoWhy} emphasizes
graph-based identification and refutation within the directed acyclic graph
(DAG) tradition championed by \citet{Pearl2009} and bridges the
potential-outcomes and graphical-model perspectives surveyed by
\citet{Imbens2020}; see \citet{SharmaKiciman2020}.
\pkg{EconML}, \pkg{DoubleML}, and \pkg{causalml} are oriented toward
machine-learning nuisance estimation, heterogeneous treatment effects, and
the double/debiased machine-learning framework of
\citet{ChernozhukovEtAl2018}
\citep{BattocchiEtAl2019,BachEtAl2022,ChenEtAl2020}. \pkg{causallib} is
closer in spirit on observational evaluation \citep{ShimoniEtAl2019}.
Design-specific tools such as \pkg{rdrobust} and \pkg{bunching} provide deep
coverage of a single design family
\citep{CalonicoEtAl2015,MavrokonstantisLockwood2020}. None of these packages
aims to carry diagnostics, sensitivity analysis, balance summaries, and
exportable review artifacts across observational, panel, IV, RD, and bunching
designs in a single workflow.

The same pattern appears in adjacent \proglang{R} software. \pkg{MatchIt}
packages matching-oriented preprocessing and balance review
\citep{HoImaiKingStuart2011}, \pkg{Synth} specializes in
synthetic-control workflows \citep{AbadieDiamondHainmueller2011}, and
\pkg{sensemakr} focuses on omitted-variable-bias sensitivity analysis
\citep{CinelliFerwerdaHazlett2020}. \citet{AbadieCattaneo2018} survey the
broader landscape of econometric methods for program evaluation, noting
that researchers routinely combine evidence across design families. Current
software typically supports that workflow only by forcing users to switch
between packages with incompatible outputs and diagnostic conventions.

CausalLens was developed to close that gap. It does not introduce a new
estimand or a new identification strategy. Its contribution is architectural:
a common result-object design that carries diagnostics, metadata, and
reproducibility artifacts alongside the effect estimate, across every
supported design family. That positioning is closer to a software-integration
contribution than to a frontier-method contribution, and it is the standard
against which the package should be judged
\citep{ImbensLemieux2008,LeeLemieux2010,AngristPischke2009,BachEtAl2022,CalonicoEtAl2015}.

\begin{table}[t!]
\centering
\footnotesize
\setlength{\tabcolsep}{4.5pt}
\begin{tabular}{L{2.4cm}L{3.2cm}L{3.2cm}L{4.2cm}}
\hline
Package & Primary focus & Design scope & Diagnostic workflow \\
\hline
\pkg{DoWhy} & Graph-based ID, refutation & Observational & Refuters, partial overlap \\
\pkg{EconML} & ML nuisance, CATE & Observational, IV & Limited built-in \\
\pkg{DoubleML} & Double/debiased ML & Observational, IV, DiD & Sensitivity via \pkg{sensemakr} \\
\pkg{causalml} & Uplift, heterogeneity & Observational & Limited built-in \\
\pkg{causallib} & Observational evaluation & Observational & Balance, overlap \\
\pkg{rdrobust} & RD bandwidth, inference & RD only & McCrary, bandwidth diagnostics \\
CausalLens & Review-oriented workflow & Obs, DiD, SC, IV, RD, Bunching & Balance, overlap, placebo, E-values, Rosenbaum bounds, OVB, export \\
\hline
\end{tabular}
\caption{\label{tab:comparison} Comparison of \proglang{Python} causal
inference packages. CausalLens is the only package that attaches a full
diagnostic and export workflow to every supported design family.}
\end{table}


%% -- Implementation ----------------------------------------------------------

\section{Package design and implementation}\label{sec:implementation}

CausalLens targets \proglang{Python}~3.11 or newer and is distributed as an
installable package with a command-line entry point, bundled benchmark data
loaders, tests, and manuscript-oriented output exports. The software is
organized around estimator classes and shared result objects so that users can
move across design families without relearning a new output format for each
method. All estimators return a \code{CausalEstimate} object that carries the
point estimate, standard error, confidence interval, $p$~value, and a nested
diagnostics container.

\subsection{Observational estimation}\label{sec:obs}

The observational core includes five estimators. Regression adjustment
estimates the average treatment effect (ATE) or average treatment effect on
the treated (ATT) via outcome modeling
\citep{Wooldridge2010}. Propensity-score matching pairs treated and control
units on estimated propensity scores following
\citet{RosenbaumRubin1983} and \citet{AbadieImbens2006}, with the matching
estimator's finite-sample properties discussed extensively by
\citet{HeckmanIchimuraTodd1997} and \citet{SmithTodd2005}.
Inverse probability weighting (IPW) reweights the sample using estimated
propensity scores as developed by \citet{HiranoImbensRidder2003}, with
practical guidance on balance assessment in IPW contexts given by
\citet{AustinStuart2015}. The standard doubly robust estimator combines an
outcome model with IPW to protect against misspecification of either
nuisance model
\citep{RobinsRotnitzkyZhao1994,BangRobins2005,FunkEtAl2011},
while the cross-fitted variant implements the sample-splitting strategy
of \citet{ChernozhukovEtAl2018} to reduce regularization bias.
\citet{KingNielsen2019} caution that propensity-score matching can increase
imbalance under certain conditions; CausalLens surfaces covariate-balance
summaries after matching so that users can detect this.

Each estimator reports reviewer-facing diagnostics. Overlap checks flag
positivity violations \citep{Stuart2010}, covariate-balance summaries report
standardized mean differences (SMD) before and after
adjustment \citep{Austin2009,AustinStuart2015}, and sensitivity
diagnostics include E-values \citep{VanderWeeleDing2017}, Rosenbaum
bounds \citep{Rosenbaum2002}, and omitted-variable-bias summaries in the
Cinelli--Hazlett framework \citep{CinelliHazlett2020}. Rather than treating
these as downstream extras, CausalLens attaches them directly to the result
object returned by every estimator.

A representative workflow for the Lalonde benchmark illustrates the
pattern:

\begin{Code}
from causal_lens import DoublyRobustEstimator, load_lalonde

data = load_lalonde()
confounders = ["age", "education", "married", "nodegree",
               "black", "hispanic", "re74", "re75"]
est = DoublyRobustEstimator("treat", "re78", confounders,
                            bootstrap_repeats = 200)
result = est.fit(data)
\end{Code}

The returned \code{result} object provides \code{result.effect},
\code{result.se}, \code{result.ci\_low}, \code{result.ci\_high},
\code{result.p\_value}, and
\code{result.diagnostics} containing overlap, balance, and sensitivity
summaries. This interface is intentionally uniform: the same attribute
names and diagnostic containers apply whether the estimator is
regression adjustment, matching, IPW, or doubly robust.

\subsection{Panel and synthetic control methods}\label{sec:panel}

For settings with longitudinal data, CausalLens provides two-period and
staggered-adoption difference-in-differences (DiD). The staggered
implementation addresses concerns raised by
\citet{GoodmanBacon2021} about heterogeneous treatment effects in
two-way fixed-effects specifications. The estimator follows the
group--time aggregation strategy of \citet{CallawaySantAnna2021}, with
doubly robust estimation drawing on \citet{SantAnnaZhao2020}. Event-study
heterogeneity considerations from \citet{SunAbraham2021} and
\citet{deChaisemartinDHaultfoeuille2020} motivate the parallel-trends
diagnostics that the package reports. Recent surveys by
\citet{RothEtAl2023} provide additional context for the design choices in
this module.

Synthetic control methods follow \citet{AbadieDiamondHainmueller2010} and
the practical guidance in \citet{Abadie2021}. Placebo-style permutation
inference is included to assess whether the estimated effect is larger than
effects obtained by reassigning treatment to donor units.

\subsection{Quasi-experimental methods}\label{sec:quasi}

Two-stage least squares (2SLS) estimation for instrumental variables
follows the identification framework of \citet{AngristImbensRubin1996}, with
the first-stage $F$~statistic reported as a weak-instrument diagnostic
following the critical values of \citet{StockYogo2005} and the rule-of-thumb
threshold of \citet{StaigerStock1997}. Sharp and fuzzy regression
discontinuity designs follow
\citet{HahnToddVanDerKlaauw2001} and the practical guidelines surveyed by
\citet{ImbensLemieux2008}, \citet{LeeLemieux2010}, and
\citet{CattaneoTitiunik2022}. Robust bias-corrected local-polynomial
inference implements the procedure of \citet{CalonicoCattaneoTitiunik2014},
and McCrary-style manipulation testing uses the density discontinuity test of
\citet{McCrary2008} with the local-polynomial density estimator of
\citet{CattaneoJanssonMa2020}. The RD module reports bandwidth, effective
sample size, and the McCrary $p$~value as standard diagnostic output.
\citet{Lee2008} provides a canonical application of sharp RD that motivates
the diagnostic defaults in the package.

\subsection{Bunching analysis}\label{sec:bunching}

The bunching module estimates excess mass at kink or notch points,
following the elasticity-oriented framework of
\citet{Saez2010} and the methodological review by \citet{Kleven2016}. The
estimator returns an excess-mass statistic, a structural elasticity
summary, and excluded-range diagnostics. Cross-design comparison functions
allow users to evaluate bunching evidence alongside observational, DiD, IV,
and RD results in a unified diagnostic table.

Table~\ref{tab:coverage} summarizes the design families currently covered
and the evidence artifacts the package returns or ships with the software.

\begin{table}[t!]
\centering
\footnotesize
\setlength{\tabcolsep}{4.5pt}
\begin{tabular}{L{3.0cm}L{5.1cm}L{5.3cm}}
\hline
Design family & Representative methods & Review-oriented outputs \\
\hline
Observational estimation & Regression adjustment, matching, IPW, DR, cross-fitted DR & Overlap checks, balance summaries, placebo tests, sensitivity diagnostics, benchmark tables \\
Panel methods & Two-period and staggered DiD & Parallel-trends review, cohort summaries, synthetic examples, cross-design diagnostics \\
Quasi-experimental methods & Synthetic control, IV, sharp and fuzzy RD & Placebo-style inference, first-stage strength, McCrary tests, exported summaries \\
Bunching analysis & Excess-mass and structural elasticity summaries & Excluded-range diagnostics, excess-mass summaries, cross-design comparison outputs \\
\hline
\end{tabular}
\caption{\label{tab:coverage} Design families covered in \pkg{CausalLens} and
the associated reviewer-facing outputs exposed in the package workflow.}
\end{table}


%% -- Validation --------------------------------------------------------------

\section{Validation and reproducible evidence}\label{sec:validation}

Validation for CausalLens rests on four complementary layers, consistent
with the emphasis on cumulative evidence that \citet{Peng2011},
\citet{SandveEtAl2013}, and \citet{StoddenEtAl2016} advocate for
computational research. The repository includes unit and integration tests
spanning estimators, diagnostics, panel methods, instrumental variables,
regression discontinuity, bunching, reporting, and reference-parity checks.
Below, the remaining validation layers are described in turn.

\subsection{Benchmark replications}\label{sec:benchmarks}

The package includes public benchmark workflows based on the Lalonde
and NHEFS datasets. The Lalonde dataset originates from
\citet{Lalonde1986}, who used it to evaluate whether nonexperimental
methods could recover the experimental benchmark, and was further analyzed by
\citet{DehejiaWahba1999}, \citet{SmithTodd2005}, and many subsequent
studies. The NHEFS dataset is the primary teaching example in
\citet{HernanRobins2020}. These benchmarks connect the software to datasets
that are already well known in the causal inference literature and allow
users to compare the package's outputs against recognizable reference values.

Table~\ref{tab:paper-benchmark-overview} provides a machine-generated summary
of the manuscript benchmark datasets generated by the same top-level
reproduction workflow used for the submission bundle.
Table~\ref{tab:paper-lalonde_public_benchmark-estimators} gives the full
Lalonde estimator-level detail, and
Table~\ref{tab:paper-nhefs_public_benchmark-estimators} provides the
corresponding NHEFS results.

\input{../../outputs/paper/tables/table01_benchmark_overview.tex}

\input{../../outputs/paper/tables/table02_lalonde_public_benchmark_estimators.tex}

\input{../../outputs/paper/tables/table03_nhefs_public_benchmark_estimators.tex}

On the Lalonde benchmark, CausalLens reports ATT estimates of \$1480.7
from regression adjustment, \$1429.1 from matching, \$1463.9 from IPW, and
\$1709.5 from the doubly robust estimator. These values fall within the
\$800--\$1800 range documented in the Lalonde replication literature
\citep{Lalonde1986,DehejiaWahba1999,SmithTodd2005,Stuart2010}.
Confidence intervals are percentile bootstrap intervals, while
$p$-values use Wald-type asymptotic inference; these two approaches may
diverge slightly in small samples.
Figure~\ref{fig:lalonde} visualizes the estimator comparison.

\begin{figure}[t!]
\centering
\includegraphics[width=0.85\textwidth]{../../outputs/paper/figures/figure01_lalonde_estimators.pdf}
\caption{\label{fig:lalonde} Lalonde benchmark: treatment-effect estimates
and 95\% confidence intervals across four observational estimators. All
estimates fall within the published \$800--\$1800 band documented in the
replication literature \citep{DehejiaWahba1999}. The figure makes clear that
the point estimates cluster in the same broad empirical range while IPW and
the doubly robust estimator retain somewhat wider uncertainty intervals than
matching and regression adjustment.}
\end{figure}

On the NHEFS benchmark, the evidence stack reports estimated weight-change
effects of 3.33~kg from regression adjustment, 3.21~kg from matching,
3.26~kg from IPW, and 3.29~kg from doubly robust estimation. These values
cluster within 0.12~kg of each other, consistent with the stability expected
when the propensity model is reasonably well specified
\citep{HernanRobins2020,FunkEtAl2011}.
Figure~\ref{fig:nhefs} visualizes the NHEFS estimator comparison, and
Figure~\ref{fig:loveplot} presents the covariate-balance love plot showing
standardized mean differences before and after propensity-score adjustment.

\begin{figure}[t!]
\centering
\includegraphics[width=0.85\textwidth]{../../outputs/paper/figures/figure02_nhefs_estimators.pdf}
\caption{\label{fig:nhefs} NHEFS benchmark: treatment-effect estimates
and 95\% confidence intervals for the smoking-cessation weight-gain
question. All estimates fall within 0.12~kg of each other
\citep{HernanRobins2020}. The close agreement across estimators emphasizes
that the empirical conclusion is not being driven by a single adjustment
strategy on this benchmark.}
\end{figure}

\begin{figure}[t!]
\centering
\includegraphics[width=0.85\textwidth]{../../outputs/paper/figures/figure03_nhefs_love_plot.pdf}
\caption{\label{fig:loveplot} NHEFS covariate-balance love plot. Standardized
mean differences before (gray) and after (black) propensity-score adjustment.
The threshold line at $|\mathrm{SMD}| = 0.1$ follows the conventional
balance criterion \citep{Austin2009,AustinStuart2015}. Most covariates move
substantially toward the threshold after adjustment, while any remaining
values above 0.1 identify the specific dimensions where residual imbalance
still deserves attention.}
\end{figure}

\subsection{Synthetic known-effect validation}\label{sec:synthetic}

The project also includes a synthetic validation dataset where the true
treatment effect is known by construction ($\tau = 2.0$). This exercise
is useful because it makes it possible to assess whether estimators and
diagnostics move coherently when the ground truth is available.
Table~\ref{tab:paper-synthetic_validation_dataset-estimators} and
Figure~\ref{fig:synthetic} summarize the results: regression adjustment
recovers $\hat{\tau} = 1.985$, IPW recovers $1.927$, doubly robust recovers
$1.935$, and matching recovers $1.707$. Regression, IPW, and doubly robust
estimates are all within 3.7\% of the true effect. Matching's larger
deviation is consistent with the finite-sample bias identified by
\citet{AbadieImbens2006} for nearest-neighbor matching without bias
correction.

\input{../../outputs/paper/tables/table04_synthetic_validation_dataset_estimators.tex}

\begin{figure}[t!]
\centering
\includegraphics[width=0.85\textwidth]{../../outputs/paper/figures/figure04_synthetic_estimators.pdf}
\caption{\label{fig:synthetic} Synthetic validation dataset: estimator
recovery of the known treatment effect ($\tau = 2.0$, dashed line).
Regression adjustment, IPW, and doubly robust estimation are all within
3.7\% of truth. The dashed reference line marks the data-generating effect,
making the downward bias of nearest-neighbor matching visually apparent
relative to the other estimators.}
\end{figure}

\subsection{Diagnostic and sensitivity validation}\label{sec:diag-val}

Two additional diagnostics are worth noting. A placebo test using
pre-treatment earnings (\code{re74}) as a fake outcome on the Lalonde
benchmark produces a non-null ``effect'' for regression adjustment, IPW,
and doubly robust estimation, while matching correctly returns a confidence
interval that includes zero (Table~\ref{tab:paper-placebo-test}). This is
the expected behavior on the Lalonde data and confirms that the diagnostic
layer surfaces residual confounding rather than hiding it
\citep{DehejiaWahba1999,Stuart2010}.

\input{../../outputs/paper/tables/table05_placebo_test.tex}

A Rosenbaum-bounds sensitivity analysis on the Lalonde matched pairs
(Table~\ref{tab:paper-rosenbaum-bounds}) shows that the matching
result is not significant even at $\Gamma = 1.0$ ($p = 0.113$) and loses
significance rapidly at higher $\Gamma$ levels, which is consistent with
the known vulnerability of the Lalonde observational comparison to hidden
bias \citep{Rosenbaum2002}.

\input{../../outputs/paper/tables/table06_rosenbaum_bounds.tex}

\subsection{Cross-design diagnostics}\label{sec:cross-design}

The cross-design diagnostic replication shows the package reporting both
reassuring and cautionary information in a single output
table. In the current synthetic comparison, the difference-in-differences
pre-trends test returns $p = 0.781$, the instrumental-variables first-stage
$F$~statistic is 144.1 (well above the \citet{StaigerStock1997}
rule-of-thumb threshold of 10 and the \citet{StockYogo2005} critical
values), the observational E-value is 4.94 \citep{VanderWeeleDing2017},
and the regression discontinuity McCrary test returns $p = 0.871$
\citep{McCrary2008,CattaneoJanssonMa2020}. At the same time, the same
summary flags a residual post-adjustment observational imbalance of
max $|\mathrm{SMD}| = 0.487$. This is precisely the kind of mixed evidence
that the package is designed to preserve rather than hide: one should be
able to export a result and its diagnostic liabilities together.

\subsection{Stability and reproducibility}\label{sec:stability}

Table~\ref{tab:paper-stability-summary} reports bootstrap stability metrics
across all three benchmark datasets and all four estimators. The
deterministic estimators exhibit zero coefficient of variation across
bootstrap resamples within each dataset, confirming numerical
reproducibility of the estimation pipeline. The balance overview in
Figure~\ref{fig:balance} provides a cross-dataset summary of covariate
balance after adjustment.

\input{../../outputs/paper/tables/table07_stability_summary.tex}

\begin{figure}[t!]
\centering
\includegraphics[width=0.85\textwidth]{../../outputs/paper/figures/figure05_balance_overview.pdf}
\caption{\label{fig:balance} Balance overview across benchmark datasets.
Mean absolute standardized mean difference after adjustment, by estimator
and dataset. Lower values indicate better covariate balance
\citep{Austin2009}. Each row summarizes how far a method moves from raw to
adjusted imbalance for a given dataset, so larger leftward shifts indicate
stronger balance improvement and the remaining endpoints show residual
post-adjustment imbalance.}
\end{figure}

Together, the unit and integration tests, benchmark replications, synthetic
known-effect recovery, diagnostic validation, and cross-design comparisons
provide a cumulative software-validation argument rather than relying on a
single benchmark or a single testing strategy. The external-comparison
audit shows exact numerical parity with hand-coded implementations across
all six checked method--dataset pairs, confirming that the estimator
internals reproduce expected results
\citep{WilsonEtAl2014,StoddenEtAl2016}. For users or reviewers evaluating
whether the software operates as claimed, the project offers several fast
checks: package tests, executable replication scripts, exported benchmark
summaries, and a command-line workflow that produces reproducible tables and
charts from bundled data.


%% -- Illustration ------------------------------------------------------------

\section[Illustration: A reviewer-facing workflow in Python]{Illustration: A reviewer-facing workflow in \proglang{Python}}\label{sec:workflow}

This section demonstrates the CausalLens workflow in three parts: the
command-line reproduction path, programmatic estimator usage, and
cross-design diagnostic comparison.

\subsection{Command-line reproduction}\label{sec:cmdline}

The package supports a reviewer-facing workflow in which a user
can install the current checkout, rerun the core evidence stack, and inspect
the generated outputs without reconstructing the manuscript logic from
scratch. The minimal command sequence is:

\begin{Code}
python -m pip install -e .
python replications/run_all.py --skip-simulation
\end{Code}

The first command installs the software from the current checkout. The second
executes the Lalonde, NHEFS, and cross-design replications and writes their
CSV outputs to \code{replications/outputs/}. When a fuller validation pass is
needed, the simulation study can also be included:

\begin{Code}
python replications/run_all.py
\end{Code}

This command reproduces all benchmark, simulation, and cross-design outputs
and regenerates the machine-generated manuscript tables and figures. The full
evidence stack completes in under five minutes on a standard workstation.

\subsection{Programmatic estimator usage}\label{sec:api}

The estimator API follows a consistent pattern across all design families.
The following example loads the NHEFS dataset and fits three observational
estimators:

\begin{Code}
from causal_lens import (RegressionAdjustmentEstimator,
    IPWEstimator, DoublyRobustEstimator, load_nhefs)

data = load_nhefs()
confounders = ["sex", "age", "race", "education",
    "smokeintensity", "smokeyrs", "exercise",
    "active", "wt71"]

for Estimator in [RegressionAdjustmentEstimator,
                  IPWEstimator, DoublyRobustEstimator]:
    est = Estimator("qsmk", "wt82_71", confounders,
                    bootstrap_repeats = 200)
    r = est.fit(data)
    print(f"{Estimator.__name__:30s}  "
          f"effect={r.effect:.3f}  "
          f"CI=[{r.ci_low:.3f}, {r.ci_high:.3f}]  "
          f"p={r.p_value:.4f}")
\end{Code}

Every returned result carries a \code{diagnostics} attribute with overlap
status, covariate-balance tables, and sensitivity statistics. A sensitivity
analysis can be requested in a single additional call:

\begin{Code}
sens = est.sensitivity_analysis(data)
print(f"E-value: {sens.e_value:.2f}  "
      f"Bias to zero: {sens.bias_to_zero_effect:.3f}")
\end{Code}

\subsection{Cross-design comparison}\label{sec:cross-api}

CausalLens exposes a unified cross-design comparison interface that
produces a single diagnostic table spanning observational, DiD, IV, RD,
and bunching designs:

\begin{Code}
from causal_lens import (DifferenceInDifferences,
    TwoStageLeastSquares, RegressionDiscontinuity,
    BunchingEstimator, compare_designs)

diagnostics = compare_designs(
    d_obs, d_did, d_iv, d_rd, d_bunch)
for d in diagnostics:
    print(f"{d.design:15s}  "
          f"{d.diagnostic_name:25s}  "
          f"value={d.diagnostic_value:.3f}  "
          f"passes={d.passes}")
\end{Code}

This interface enables the kind of multi-design evidence review advocated
by \citet{AbadieCattaneo2018} and \citet{AtheyImbens2017}, where the
credibility of a causal claim is strengthened by convergent evidence across
designs with different identifying assumptions.


%% -- Discussion --------------------------------------------------------------

\section{Discussion}\label{sec:discussion}

CausalLens occupies a specific niche: it is most useful where analysts value a
consistent, auditable workflow across design families more than frontier
performance within any single one. The design deliberately trades
specialization depth for cross-design consistency, following the
observation by \citet{AbadieCattaneo2018} that credible causal inference
often involves comparing estimates from multiple designs with different
identifying assumptions.

The package does not replace graph-based
identification frameworks \citep{Pearl2009,SharmaKiciman2020},
machine-learning-first treatment-effect platforms
\citep{BattocchiEtAl2019,BachEtAl2022,ChenEtAl2020},
state-of-the-art RD toolchains
\citep{CalonicoEtAl2015,CattaneoIdroboTitiunik2020,CattaneoTitiunik2022},
matching-oriented preprocessing packages \citep{HoImaiKingStuart2011},
synthetic-control software \citep{AbadieDiamondHainmueller2011},
omitted-variable-bias sensitivity packages
\citep{CinelliFerwerdaHazlett2020}, or specialized bunching implementations
\citep{MavrokonstantisLockwood2020}. Its contribution is that estimates
from several established designs arrive with their diagnostics instead of
leaving those checks to a separate notebook or reporting layer.

The validation evidence supports that claim concretely. Lalonde benchmark
estimates fall in the published \$800--\$1800 band
\citep{Lalonde1986,DehejiaWahba1999}, NHEFS estimates cluster within
0.12~kg of each other \citep{HernanRobins2020}, synthetic known-effect
recovery is within 3.7\% of truth for regression, IPW, and doubly robust
estimation \citep{FunkEtAl2011}, and the external-comparison audit shows
exact numerical parity with hand-coded implementations across all six
checked method--dataset pairs. Placebo tests and Rosenbaum bounds correctly
surface the known limitations of matching on the Lalonde and NHEFS
benchmarks \citep{Rosenbaum2002,VanderWeeleDing2017}, demonstrating that
the diagnostic layer flags problems rather than concealing them.

\subsection{Limitations and scope boundaries}\label{sec:limitations}

Several limitations are worth noting. The propensity-score models use
logistic regression with \pkg{scikit-learn}'s default solver; users who need
flexible nonparametric models (e.g., random forests, gradient boosting) for
nuisance estimation should consider \pkg{EconML} or \pkg{DoubleML}
\citep{BattocchiEtAl2019,BachEtAl2022}. The matching estimator uses
nearest-neighbor matching without the bias correction of
\citet{AbadieImbens2006}, which contributes to the finite-sample bias
visible in Table~\ref{tab:paper-synthetic_validation_dataset-estimators}.
The RD module does not yet implement the optimal bandwidth selection of
\citet{CalonicoCattaneoFarrell2020} or the coverage-error-rate-optimal
procedures discussed by \citet{CattaneoTitiunik2022}. The DiD module does
not yet support the full staggered-adoption heterogeneity diagnostics
surveyed by \citet{RothEtAl2023}.

\subsection{Future work}\label{sec:future}

Natural extensions include graph-based identification and refutation
\citep{Pearl2009,Imbens2020}, heterogeneous treatment-effect estimation
via machine-learning nuisance models
\citep{ChernozhukovEtAl2018,BattocchiEtAl2019}, advanced RD bandwidth
selection \citep{CalonicoCattaneoFarrell2020,ImbensKalyanaraman2012},
notch-specific bunching extensions \citep{Kleven2016}, and publication on
the Python Package Index (PyPI) to satisfy JSS's preference for standard
repository availability. These extensions do not change the current
contribution: a compact, testable, replication-oriented \proglang{Python}
package for auditable causal analysis that keeps diagnostic context attached
to every estimate it produces.


%% -- Computational details ---------------------------------------------------

\section*{Computational details}

The current release (version 0.5.0) targets \proglang{Python}~3.11 or newer
on Windows, macOS, and Linux. Core dependencies include
\pkg{NumPy} $\geq 1.26$ \citep{NumPyEtAl2020},
\pkg{pandas} $\geq 2.2$ \citep{McKinney2010},
\pkg{scikit-learn} $\geq 1.5$ \citep{PedregosaEtAl2011},
\pkg{SciPy} $\geq 1.13$ \citep{VirtanenEtAl2020},
\pkg{statsmodels} $\geq 0.14$ \citep{SeaboldPerktold2010}, and
\pkg{Matplotlib} $\geq 3.9$ \citep{Hunter2007}. Package metadata and dependency constraints are
declared in \code{pyproject.toml}. The implementation relies on
widely used scientific \proglang{Python} libraries, especially
\pkg{scikit-learn} for nuisance modeling and \pkg{statsmodels} for
regression-based estimation and inference.
Replication scripts write outputs to
the repository's \code{replications/outputs/} directory, while manuscript-
oriented exports are written under \code{outputs/paper/}.
The complete evidence stack (Lalonde, NHEFS, cross-design, and simulation
replications plus manuscript table and figure generation) completes in under
five minutes on a standard workstation. The replication log and generated
outputs are included in the submission bundle to facilitate review.


\section*{Acknowledgments}

CausalLens was developed as an independent research project. No external
funding was received. The author declares no competing interests. The author
thanks the open-source contributors to
\pkg{statsmodels}, \pkg{scikit-learn}, \pkg{pandas}, \pkg{NumPy}, and
\pkg{SciPy}, on which the package depends.


\bibliography{causal-lens}


\end{document}
